\section{受案范围}
\begin{enumerate}
  \item 每年都有。
  \item \zd{具体行政行为}可受。区分\qd{司法行为、抽象行政行为、行政内部行为}等。\al{征求意见，复函}。
  \item \zd{注销撤回撤销}是一种独立的可诉具体性的行为。
  \item 信访行为\zd{不能}诉。提防\qd{假信访}。
  \item \zd{行政协议}虽然不是具体行政行为，但\zd{可受}。
  \item \zd{部分}的抽象行为可以\zd{附带性}（不是\qd{直接}）的受案。必须和具体行为之间\zd{有关联性}。
        \begin{enumerate}
          \item 部分：\zd{行政法规、规章}（省市政府、国务院部门）、\zd{国务院制定}的除外。
          \item 原则\zd{开庭前}提出，例外：有正当理由可以\zd{法庭调查中}提出。
          \item \zd{具体行为}来决定管辖法院。
          \item 审查程序和实体要求。
          \item 听取意见（前提\zd{认为文件可能不合法}）。审查过程当中应当听取意见。错。很隐蔽。
          \item 审查结果：
                \begin{enumerate}
                  \item 法院不能宣告文件无效 \hz 撤销（\zd{没有}违宪违法审查权）。
                  \item \zd{做出裁判}的法院，裁判\zd{生效}之后的\zd{3 个月内}向\zd{制定机关}提出修改 \hz 废止的建议。\qd{和复议有区别}。
                  \item 不适用提建议告领导（县政府市公安检察机关）。
                  \item 裁判生效后\zd{分别层报}：\qd{省政府 \yjt 省高院；国务院 \yjt 最高法}。
                \end{enumerate}
        \end{enumerate}

        每年一题：制了一个抽象行为，如一个工作方案/通知，回头做了一个具体行为，做了一个处罚（责令关闭之类）。
\end{enumerate}
